
Thus let @@M0@@ be a torsor under @@M1@@, corresponding to an extension @@M2@@ of @@M3@@ by @@M4@@. The algebra @@M5@@ of @@M6@@ over @@M7@@ is @@M8@@, that is, @@M9@@, which identifies with the ring of homogeneous fractions of degree @@M10@@, @@M11@@. Consequently @@M12@@ identifies with the open @@M13@@ of @@M14@@. The complement of @@M15@@ in @@M16@@ is then @@M17@@, which identifies, by (3.3.1), with @@M18@@. Thus one has the immersions:
@@M19@@
where @@M20@@ is closed in @@M21@@, and @@M22@@ is open and complementary to @@M23@@.

\sgaStableTarget{sga6.e09.prop.3-4}\textbf{Proposition 3.4.} \uline{Let} @@M24@@ \uline{be a noetherian scheme}, @@M25@@ \uline{a locally free} @@M26@@\uline{-Module of rank} @@M27@@, @@M28@@ \uline{the associated vector bundle}, @@M29@@ \uline{a torsor under} @@M30@@, \uline{and} @@M31@@ \uline{the structural morphism of} @@M32@@. \uline{Then the canonical homomorphism}:
@@M33@@
\uline{is an isomorphism}.

With the identifications of 3.3, one has the exact sequence:
\sgaStableTarget{sga6.e09.eq.3-4-1}@@M34@@
where @@M35@@ denotes the closed immersion @@M36@@ defined by @@M37@@. By 3.2, one has @@M38@@, @@M39@@.
Since @@M40@@ and @@M41@@ are smooth, @@M42@@ is a regular immersion, and one therefore has a homomorphism @@M43@@; the homomorphism between the @@M44@@ is then obtained by tensoring the homomorphism between the @@M45@@ by @@M46@@: to see this, it is enough to show that for @@M47@@ and

